% !TEX root = ../main.tex
% Tóm tắt nghiên cứu: bối cảnh, động lực, phương pháp và đóng góp.
% Biên dịch độc lập qua workflow; không thay đổi định dạng thesis.
\frontpage{TÓM TẮT}

Trong cộng tác người--robot, chất lượng phối hợp không chỉ phụ thuộc vào việc robot thực hiện đúng thao tác mà còn ở thời điểm và vị trí hành động so với tiến độ của con người. Nếu chờ đến khi hành động của người hoàn tất, robot có thể bỏ lỡ cơ hội làm việc đồng thời. Ngược lại, hành động quá sớm dựa trên nhận định sai có thể dẫn đến thao tác không phù hợp hoặc phải thực hiện lại. Vì vậy, robot cần lựa chọn và điều chỉnh hành động khi mục tiêu của người chưa được quan sát đầy đủ và có thể thay đổi.

Khóa luận nghiên cứu cách học chính sách hành động cộng tác từ trình diễn, kết hợp quan sát hiện tại với dự đoán chuyển động tương lai của người. Động lực chính là dữ liệu trình diễn có thể thể hiện cách chuyên gia phối hợp, nhưng việc bổ sung dự đoán vào đầu vào chưa bảo đảm chính sách học được khi nào cần sử dụng thông tin đó. Nghiên cứu sử dụng Action Chunking with Transformers (ACT) để học hành động robot và ManiCast để cung cấp quỹ đạo chuyển động dự đoán. Trọng tâm phương pháp là biểu diễn cấu trúc không gian và thời gian của quỹ đạo, tăng chú trọng vào các thời điểm quyết định trong huấn luyện và duy trì khả năng cập nhật hành động khi quan sát mới khác với dự đoán.

Nghiên cứu xây dựng các nhiệm vụ cộng tác trong môi trường mô phỏng có ràng buộc vật lý, tạo trình diễn chuyên gia và đánh giá chính sách trên những tình huống thông thường, dự đoán sai và thay đổi ý định. Đóng góp của khóa luận nằm ở việc xây dựng và khảo sát một phương pháp học chính sách hành động thống nhất có điều kiện bởi dự đoán chuyển động, cùng quy trình đánh giá tách biệt khả năng hoàn thành nhiệm vụ, hành động sớm, phối hợp thời gian và xử lý sai lệch dự đoán. Qua đó, khóa luận làm rõ điều kiện để thông tin về chuyển động tương lai hỗ trợ quyết định hành động của robot, thay vì chỉ xuất hiện như một đầu vào bổ sung.
